\chapter{2012年天津卷（文）}

\section{单选题}

\begin{problem}
\(\i\) 为虚数单位，复数 \(\frac{5 + 3\i}{4 - \i} =\)（\quad）

\choices
    {\(1 - \i\)}
    {\(-1 + \i\)}
    {\(1 + \i\)}
    {\(-1 - \i\)}
\end{problem}

\begin{answer}
C
\end{answer}

\begin{solution}
分子分母同乘以分母的共轭复数，得 \(\frac{5+3\i}{4-\i}=\frac{(5+3\i)(4+\i)}{(4-\i)(4+\i)}=\frac{20+5\i+12\i+3\i^2}{17}=\frac{17+17\i}{17}=1+\i\). 因此复数对应选项 C.
\end{solution}

\begin{problem}
设变量 \(x\)，\(y\) 满足约束条件 \(\begin{cases}
2x + y - 2 \geqslant 0,\\
x - 2y + 4 \geqslant 0,\\
x - 1 \leqslant 0,
\end{cases}\) 则目标函数 \(z = 3x - 2y\) 的最小值为（\quad）

\choices
    {\(-5\)}
    {\(-4\)}
    {\(-2\)}
    {3}
\end{problem}

\begin{answer}
B
\end{answer}

\begin{solution}
约束条件等价于 \(y\geqslant2-2x\)，\(y\leqslant\frac{x}{2}+2\)，\(x\leqslant1\). 由上下界能够同时满足，得 \(2-2x\leqslant\frac{x}{2}+2\)，所以 \(x\geqslant0\). 又有 \(x\leqslant1\)，故 \(0\leqslant x\leqslant1\). 对固定的 \(x\)，由于 \(z=3x-2y\) 随 \(y\) 增大而减小，取 \(y=\frac{x}{2}+2\) 时取得最小值下界，于是 \(z\geqslant3x-2\left(\frac{x}{2}+2\right)=2x-4\geqslant-4\). 当 \((x,y)=(0,2)\) 时满足全部约束且 \(z=-4\)，所以最小值为 \(-4\)，故选 B.
\end{solution}

\begin{problem}
阅读如图程序框图，运行相应的程序，则输出 \(S\) 的值为（\quad）

\bitmapfigure[max width=0.82\linewidth]{img/2012/tianjin_liberal/q3_fig1.png}

\choices
    {8}
    {18}
    {26}
    {80}
\end{problem}

\begin{answer}
C
\end{answer}

\begin{solution}
程序开始时 \(n=1\)，\(S=0\). 第一次循环后 \(S=0+3^1-3^0=2\)，\(n=2\)；第二次循环后 \(S=2+3^2-3^1=8\)，\(n=3\)；第三次循环后 \(S=8+3^3-3^2=26\)，\(n=4\). 此时满足 \(n\geqslant4\)，输出 \(S=26\)，故选 C.
\end{solution}

\begin{problem}
已知 \(a = 2^{1.2}\)，\(b = \left(\frac{1}{2}\right)^{-0.8}\)，\(c = 2\log_5 2\)，则 \(a\)，\(b\)，\(c\) 的大小关系为（\quad）

\choices
    {\(c < b < a\)}
    {\(c < a < b\)}
    {\(b < a < c\)}
    {\(b < c < a\)}
\end{problem}

\begin{answer}
A
\end{answer}

\begin{solution}
因为 \(b=\left(\frac12\right)^{-0.8}=2^{0.8}=2^{\frac45}\)，而 \(a=2^{1.2}=2^{\frac65}\)，底数 \(2>1\)，所以 \(b<a\). 又因 \(2<\sqrt5\)，有 \(\log_5 2<\frac12\)，从而 \(c=2\log_5 2<1\). 同时 \(b=2^{\frac45}>1\)，故 \(c<b<a\)，选项为 A.
\end{solution}

\begin{problem}
设 \(x \in \R\)，则“\(x > \frac{1}{2}\)”是“\(2x^2 + x - 1 > 0\)”的（\quad）

\choices
    {充分而不必要条件}
    {必要而不充分条件}
    {充分必要条件}
    {既不充分也不必要条件}
\end{problem}

\begin{answer}
A
\end{answer}

\begin{solution}
将二次式因式分解，得 \(2x^2+x-1=(2x-1)(x+1)\). 因此 \(2x^2+x-1>0\) 的解集为 \(x<-1\) 或 \(x>\frac12\). 条件 \(x>\frac12\) 能推出不等式成立，但不等式成立还允许 \(x<-1\)，所以前者是后者的充分而不必要条件，故选 A.
\end{solution}

\begin{problem}
下列函数中，既是偶函数，又在区间 \((1,2)\) 内是增函数的函数为（\quad）

\choices
    {\(y = \cos 2x, x \in \R\)}
    {\(y = \log_2 |x|\)，\(x \in \R\) 且 \( x \neq 0 \)}
    {\(y = \frac{\e^x - \e^{-x}}{2}\)，\(x \in \R\)}
    {\(y = x^{3} + 1\)，\(x \in \R\)}
\end{problem}

\begin{answer}
B
\end{answer}

\begin{solution}
选项 A 中，\(\cos2x\) 是偶函数，但其导数 \(-2\sin2x\) 在 \((1,2)\) 内变号，故不是增函数. 选项 B 中，定义域为 \(\R\mathbin{\backslash}\{0\}\)，且 \(\log_2|-x|=\log_2|x|\)，所以是偶函数；当 \(x\in(1,2)\) 时，导数 \(\frac{1}{x\ln2}>0\)，所以在该区间内是增函数. 选项 C 满足 \(f(-x)=-f(x)\)，是奇函数；选项 D 有 \(f(-x)=-x^3+1\ne x^3+1\)，不是偶函数. 因此选 B.
\end{solution}

\begin{problem}
将函数 \( f(x) = \sin \omega x \) （其中 \( \omega > 0 \)） 的图象向右平移 \( \frac{\pi}{4} \) 个单位长度，所得图象经过点 \( \left(\frac{3\pi}{4}, 0\right) \)，则 \( \omega \) 的最小值是（\quad）

\choices
    {\( \frac{1}{3} \)}
    {1}
    {\( \frac{5}{3} \)}
    {2}
\end{problem}

\begin{answer}
D
\end{answer}

\begin{solution}
向右平移 \(\frac{\pi}{4}\) 后的函数为 \(y=\sin\left[\omega\left(x-\frac{\pi}{4}\right)\right]\). 图象经过 \(\left(\frac{3\pi}{4},0\right)\)，所以 \(\sin\frac{\omega\pi}{2}=0\). 因而 \(\frac{\omega\pi}{2}=k\pi\)，其中 \(k\in\Z\). 由 \(\omega>0\) 得 \(k\) 为正整数，故 \(\omega=2k\) 的最小值为 \(2\)，选 D.
\end{solution}

\begin{problem}
在 \(\triangle ABC\) 中，\(\angle A = 90^{\circ}\)，\(AB = 1\)，\(AC = 2\). 设点 \(P\)，\(Q\) 满足 \(\overrightarrow{AP} = \lambda \overrightarrow{AB}\)，\(\overrightarrow{AQ} = (1 - \lambda)\overrightarrow{AC}\)，\(\lambda \in \R\). 若 \(\overrightarrow{BQ} \cdot \overrightarrow{CP} = -2\)，则 \(\lambda =\)（\quad）

\choices
    {\( \frac{1}{3} \)}
    {\( \frac{2}{3} \)}
    {\( \frac{4}{3} \)}
    {2}
\end{problem}

\begin{answer}
B
\end{answer}

\begin{solution}
建立直角坐标系，令 \(A(0,0)\)，\(B(1,0)\)，\(C(0,2)\). 由向量条件得 \(P(\lambda,0)\)，\(Q(0,2-2\lambda)\)，所以 \(\overrightarrow{BQ}=(-1,2-2\lambda)\)，\(\overrightarrow{CP}=(\lambda,-2)\). 于是 \(\overrightarrow{BQ}\cdot\overrightarrow{CP}=-\lambda-2(2-2\lambda)=3\lambda-4\). 由题设 \(3\lambda-4=-2\)，解得 \(\lambda=\frac23\)，故选 B.
\end{solution}

\section{填空题}

\begin{problem}
集合 \(A = \{x \in \R \mid |x - 2| \leqslant 5\}\) 中的最小整数为\fillinblank{}.
\end{problem}

\begin{answer}
\(-3\)
\end{answer}

\begin{solution}
由 \(\lvert x-2\rvert\leqslant5\)，得 \(-5\leqslant x-2\leqslant5\)，即 \(-3\leqslant x\leqslant7\). 因此集合中的最小整数为 \(-3\).
\end{solution}

\begin{problem}
一个几何体的三视图如图所示（单位：\(\mathrm{m}\)），则该几何体的体积为\fillinblank{}\(\mathrm{m}^3\).

\bitmapfigure[max width=0.82\linewidth]{img/2012/tianjin_liberal/q10_fig1.png}
\end{problem}

\begin{answer}
30
\end{answer}

\begin{solution}
由正视图和侧视图可知，几何体的底部是长、宽、高分别为 \(3\)，\(4\)，\(2\) 的长方体，体积为 \(3\times4\times2\). 由正视图可知，底部上方的截面由一个边长为 \(1\)、高为 \(1\) 的直角三角形和一个长、宽分别为 \(1\)、\(1\) 的矩形组成；由俯视图可知该截面沿长度为 \(4\) 的方向延伸，因此上方部分的体积为 \(\left(\frac12\times1\times1+1\times1\right)\times4\). 所以几何体的体积为 \(3\times4\times2+\left(\frac12\times1\times1+1\times1\right)\times4=24+6=30\ \mathrm{m}^3\).
\end{solution}

\begin{problem}
已知双曲线 \(C_{1}:\frac{x^{2}}{a^{2}} -\frac{y^{2}}{b^{2}} = 1\) （\(a > 0\)，\(b > 0\)）与双曲线 \(C_2:\frac{x^2}{4} -\frac{y^2}{16} = 1\) 有相同的渐近线，且 \(C_1\) 的右焦点为 \(F(\sqrt{5},0)\)，则 \(a =\fillinblank{}\)，\(\ b =\fillinblank{}\).
\end{problem}

\begin{answer}
\(1,2\)
\end{answer}

\begin{solution}
双曲线 \(C_2\) 的渐近线为 \(y=\pm2x\)，而 \(C_1\) 的渐近线为 \(y=\pm\frac{b}{a}x\)，所以 \(\frac{b}{a}=2\)，即 \(b=2a\). 设 \(C_1\) 的焦半距为 \(c_1\)，则 \(c_1^2=a^2+b^2=5a^2\). 由右焦点为 \(F(\sqrt5,0)\)，得 \(c_1=\sqrt5\)，所以 \(a=1\)，进而 \(b=2\).
\end{solution}

\begin{problem}
设 \(m\)，\(n \in \R\)，若直线 \(l: mx + ny - 1 = 0\) 与 \(x\) 轴相交于点 \(A\)，与 \(y\) 轴相交于点 \(B\)，且 \(l\) 与圆 \(x^2 + y^2 = 4\) 相交所得弦的长为 \(2\)，\(O\) 为坐标原点，则 \(\triangle AOB\) 面积的最小值为\fillinblank{}.
\end{problem}

\begin{answer}
3
\end{answer}

\begin{solution}
因为直线与两坐标轴相交，故 \(m\ne0\)，\(n\ne0\)，且其到原点的距离为 \(d=\frac{1}{\sqrt{m^2+n^2}}\). 圆的半径为 \(2\)，弦长为 \(2\)，所以半弦长为 \(1\)，由直角三角形得 \(d^2+1^2=2^2\)，即 \(d=\sqrt3\). 从而 \(m^2+n^2=\frac13\).

直线在坐标轴上的截距分别为 \(\frac1m\) 和 \(\frac1n\)，故 \(\triangle AOB\) 的面积为 \(S=\frac{1}{2\lvert mn\rvert}\). 由 \(m^2+n^2\geqslant2\lvert mn\rvert\)，得 \(\lvert mn\rvert\leqslant\frac16\)，所以 \(S\geqslant3\). 当 \(\lvert m\rvert=\lvert n\rvert=\frac{1}{\sqrt6}\) 时等号成立，故面积的最小值为 \(3\).
\end{solution}

\begin{problem}
如图，已知 \(AB\) 和 \(AC\) 是圆的两条弦，过点 \(B\) 作圆的切线与 \(AC\) 的延长线相交于点 \(D\). 过点 \(C\) 作 \(BD\) 的平行线与圆相交于点 \(E\)，与 \(AB\) 相交于点 \(F\)，\(AF = 3\)，\(FB = 1\)，\(EF = \frac{3}{2}\)，则线段 \(CD\) 的长为\fillinblank{}.
\bitmapfigure[max width=0.66\linewidth]{img/2012/tianjin_liberal/q13_fig1.png}
\end{problem}

\begin{answer}
\(\frac{4}{3}\)
\end{answer}

\begin{solution}
由相交弦定理，得 \(AF\cdot FB=FC\cdot FE\)，所以 \(3\times1=FC\times\frac32\)，从而 \(FC=2\).

因为 \(CE\parallel BD\)，所以 \(\triangle AFC\sim\triangle ABD\)，于是 \(\frac{AF}{AB}=\frac{AC}{AD}=\frac{FC}{BD}=\frac34\). 设 \(CD=x\)，则 \(AD=4x\)，\(AC=3x\). 由点 \(D\) 的切割线定理，\(DB^2=DC\cdot DA=4x^2\)，故 \(DB=2x\). 又由相似比 \(\frac{FC}{BD}=\frac34\)，得 \(\frac{2}{2x}=\frac34\)，所以 \(x=\frac43\). 因此 \(CD=\frac43\).
\end{solution}

\begin{problem}
已知函数 \( y = \frac{|x^2 - 1|}{x - 1} \) 的图象与函数 \( y = kx \) 的图象恰有两个交点，则实数 \( k \) 的取值范围是\fillinblank{}.
\end{problem}

\begin{answer}
\((0,1) \cup (1,2)\)
\end{answer}

\begin{solution}
因 \(x\ne1\)，将函数分段化简为 \(y=x+1\)（\(x<-1\) 或 \(x>1\)），以及 \(y=-x-1\)（\(-1\leqslant x<1\)）. 与直线 \(y=kx\) 联立时，第一、三段都给出 \(x=\frac{1}{k-1}\)，第二段给出 \(x=-\frac{1}{k+1}\)，并分别检验定义域和分段范围.

当 \(x=\frac{1}{k-1}<-1\) 时，必须有 \(0<k<1\)；当 \(x=\frac{1}{k-1}>1\) 时，必须有 \(1<k<2\). 因而外侧两段合计在 \(0<k<1\) 或 \(1<k<2\) 时各产生一个交点，其余 \(k\) 值不产生外侧交点（\(k=1\) 时方程无解）.

对中间一段，\(-1\leqslant-\frac{1}{k+1}<1\) 的解为 \(k\in(-\infty,-2)\cup[0,+\infty)\)，其中 \(k=-2\) 对应 \(x=1\)，因而不在该段定义域内；\(k=0\) 对应中间段唯一交点 \((-1,0)\). 所以 \(k<-2\) 或 \(k\geqslant0\) 时中间段各有一个交点，其他情形没有中间段交点.

因此，恰有两个交点的情形正好是外侧一交点且中间一交点，即 \(k\in(0,1)\cup(1,2)\).
\end{solution}

\section{解答题}

\begin{problem}
某地区有小学 \(21\text{ 所}\)，中学 \(14\text{ 所}\)，大学 \(7\text{ 所}\)，现采用分层抽样的方法从这些学校中抽取 \(6\text{ 所}\) 学校对学生进行视力调查.
\begin{enumerate}
\item 求应从小学，中学，大学中分别抽取的学校数目；
\item 若从抽取的 \(6\text{ 所}\) 学校中随机抽取 \(2\text{ 所}\) 学校做进一步数据分析，
\begin{enumerate}
\item 列出所有可能的抽取结果；
\item 求抽取的 \(2\text{ 所}\) 学校均为小学的概率.
\end{enumerate}
\end{enumerate}
\end{problem}

\begin{solution}
\begin{enumerate}
\item 全部学校数为 \(21+14+7=42\)，抽样比例为 \(\frac{6}{42}=\frac17\). 因此应从小学、中学、大学中分别抽取 \(21\times\frac17=3\)，\(14\times\frac17=2\)，\(7\times\frac17=1\) 所.
\item
\begin{enumerate}
\item 设抽取的小学为 \(A_1\)，\(A_2\)，\(A_3\)，中学为 \(B_1\)，\(B_2\)，大学为 \(C_1\). 从这 \(6\) 所学校中抽取 \(2\) 所的所有可能结果为 \(\{A_1,A_2\}\)，\(\{A_1,A_3\}\)，\(\{A_1,B_1\}\)，\(\{A_1,B_2\}\)，\(\{A_1,C_1\}\)，\(\{A_2,A_3\}\)，\(\{A_2,B_1\}\)，\(\{A_2,B_2\}\)，\(\{A_2,C_1\}\)，\(\{A_3,B_1\}\)，\(\{A_3,B_2\}\)，\(\{A_3,C_1\}\)，\(\{B_1,B_2\}\)，\(\{B_1,C_1\}\)，\(\{B_2,C_1\}\)，共 \(15\) 种.
\item 其中两所均为小学的结果有 \(\{A_1,A_2\}\)，\(\{A_1,A_3\}\)，\(\{A_2,A_3\}\)，共 \(3\) 种. 所以所求概率为 \(\frac{3}{15}=\frac15\).
\end{enumerate}
\end{enumerate}
\end{solution}

\begin{problem}
在 \(\triangle ABC\) 中，内角 \(A\)，\(B\)，\(C\) 所对的边分别是 \(a\)，\(b\)，\(c\). 已知 \(a = 2\)，\(c = \sqrt{2}\)，\(\cos A = -\frac{\sqrt{2}}{4}\).
\begin{enumerate}
\item 求 \(\sin C\) 和 \(b\) 的值；
\item 求 \( \cos\left(2A+\frac{\pi}{3}\right) \) 的值.
\end{enumerate}
\end{problem}

\begin{solution}
\begin{enumerate}
\item 因为 \(A\) 是三角形内角，所以 \(\sin A=\sqrt{1-\cos^2A}=\sqrt{1-\frac18}=\frac{\sqrt{14}}4\). 由正弦定理，\(\frac{a}{\sin A}=\frac{c}{\sin C}\)，所以 \(\sin C=\frac{c\sin A}{a}=\frac{\sqrt2\cdot\frac{\sqrt{14}}4}{2}=\frac{\sqrt7}{4}\).

由余弦定理，\(a^2=b^2+c^2-2bc\cos A\)，代入已知量得 \(4=b^2+2-2b\sqrt2\left(-\frac{\sqrt2}{4}\right)=b^2+b+2\). 因而 \(b^2+b-2=0\)，即 \((b-1)(b+2)=0\). 由于 \(b>0\)，所以 \(b=1\).
\item 由二倍角公式，\(\cos2A=2\cos^2A-1=-\frac34\)，\(\sin2A=2\sin A\cos A=2\cdot\frac{\sqrt{14}}4\cdot\left(-\frac{\sqrt2}{4}\right)=-\frac{\sqrt7}{4}\). 因此 \(\cos\left(2A+\frac{\pi}{3}\right)=\cos2A\cos\frac{\pi}{3}-\sin2A\sin\frac{\pi}{3}=-\frac34\cdot\frac12+\frac{\sqrt7}{4}\cdot\frac{\sqrt3}{2}=\frac{-3+\sqrt{21}}8\).
\end{enumerate}
\end{solution}

\begin{problem}
如图，在四棱锥 \(P - ABCD\) 中，底面 \(ABCD\) 是矩形，\(AD \perp PD\)，\(BC = 1\)，\(PC = 2\sqrt{3}\)，\(PD = CD = 2\).
\begin{enumerate}
\item 求异面直线 \(PA\) 与 \(BC\) 所成角的正切值；
\item 证明：平面 \(PDC \perp\) 平面 \(ABCD\)；
\item 求直线 \(PB\) 与平面 \(ABCD\) 所成角的正弦值.
\end{enumerate}
\bitmapfigure[max width=0.60\linewidth]{img/2012/tianjin_liberal/q17_fig1.png}
\end{problem}

\begin{solution}
\begin{enumerate}
    \item 建立空间直角坐标系，使底面 \(ABCD\) 在 \(z=0\) 平面内，取 \(A(0,0,0)\)，\(B(2,0,0)\)，\(C(2,1,0)\)，\(D(0,1,0)\). 设 \(P(x,y,z)\). 由 \(AD\perp PD\) 得 \(y=1\)，由 \(PD=2\) 得 \(x^2+z^2=4\). 又由 \(PC=2\sqrt3\) 得 \((x-2)^2+z^2=12\). 两式相减得 \(x=-1\)，从而 \(z^2=3\). 取坐标系的正侧使 \(z>0\)，则 \(z=\sqrt3\)，即 \(P(-1,1,\sqrt3)\)；若取另一侧，只是关于底面的对称，以下长度和角度不变.

直线 \(PA\) 和 \(BC\) 的方向向量可分别取 \(\overrightarrow{AP}=(-1,1,\sqrt3)\) 和 \(\overrightarrow{BC}=(0,1,0)\). 设它们所成的锐角为 \(\theta\)，则 \(\cos\theta=\frac{1}{\sqrt5}\)，所以 \(\sin\theta=\frac{2}{\sqrt5}\)，从而 \(\tan\theta=2\).

\item 在矩形 \(ABCD\) 中，\(AD\perp DC\). 题设又有 \(AD\perp PD\)，而 \(DC\) 与 \(PD\) 是平面 \(PDC\) 内相交的两条直线，所以由直线与平面垂直的判定定理，得 \(AD\perp\text{平面 }PDC\). 因为 \(AD\) 在平面 \(ABCD\) 内，所以平面 \(PDC\perp\) 平面 \(ABCD\).

\item 点 \(P\) 在底面的射影为 \(H(-1,1,0)\)，而 \(B(2,0,0)\)，故 \(BH=\sqrt{(-3)^2+1^2}=\sqrt{10}\)，\(PH=\sqrt3\)，\(PB=\sqrt{BH^2+PH^2}=\sqrt{13}\). 设直线 \(PB\) 与平面 \(ABCD\) 所成的角为 \(\alpha\)，则 \(\sin\alpha=\frac{PH}{PB}=\frac{\sqrt3}{\sqrt{13}}=\frac{\sqrt{39}}{13}\).
\end{enumerate}
\end{solution}

\begin{problem}
已知 \(\{a_{n}\}\) 是等差数列，其前 \(n\) 项和为 \(S_{n}\)，\(\{b_{n}\}\) 是等比数列，且 \(a_{1} = b_{1} = 2\)，\(a_{4} + b_{4} = 27\)，\(S_{4} - b_{4} = 10\).
\begin{enumerate}
\item 求数列 \( \{a_{n}\} \) 与 \( \{b_{n}\} \) 的通项公式；
\item 记 \(T_{n} = a_{1}b_{1} + a_{2}b_{2} + \dots + a_{n}b_{n}\)，\(n \in \mathbf{N}^{*}\)，证明： \(T_{n} - 8 = a_{n-1}b_{n+1}\)（\(n \in \mathbf{N}^{*}\)，\(n > 2\)）.
\end{enumerate}
\end{problem}

\begin{solution}
\begin{enumerate}
\item 设等差数列的公差为 \(d\)，等比数列的公比为 \(q\). 则 \(a_4=2+3d\)，\(S_4=2(a_1+a_4)=4+2a_4\)，\(b_4=2q^3\). 由已知条件得 \(a_4+b_4=27\)，以及 \(4+2a_4-b_4=10\). 联立两式，得 \(a_4=11\)，\(b_4=16\). 所以 \(d=\frac{11-2}{3}=3\)，且 \(2q^3=16\)，即 \(q=2\). 因此 \(a_n=2+3(n-1)=3n-1\)，\(b_n=2\cdot2^{n-1}=2^n\).
\item 由通项公式，对 \(k\geqslant3\) 有
\(a_kb_k=(3k-1)2^k=(3k-4)2^{k+1}-(3k-7)2^k=a_{k-1}b_{k+1}-a_{k-2}b_k\).
对 \(n>2\) 求和，得
\(\displaystyle\sum_{k=3}^{n}a_kb_k=a_{n-1}b_{n+1}-a_1b_3\).
又 \(a_1b_1=4\)，\(a_2b_2=5\times4=20\)，\(a_1b_3=2\times8=16\)，所以
\(T_n=4+20+a_{n-1}b_{n+1}-16=a_{n-1}b_{n+1}+8\).
因此 \(T_n-8=a_{n-1}b_{n+1}\).
\end{enumerate}
\end{solution}

\begin{problem}
已知椭圆 \(\frac{x^2}{a^2} + \frac{y^2}{b^2} = 1\)（\(a > b > 0\)），点 \(P\left(\frac{\sqrt{5}}{5}a, \frac{\sqrt{2}}{2}a\right)\) 在椭圆上.
\begin{enumerate}
\item 求椭圆的离心率；
\item 设 \(A\) 为椭圆的左顶点，\(O\) 为坐标原点. 若点 \(Q\) 在椭圆上且满足 \(|AQ| = |AO|\)，求直线 \(OQ\) 的斜率的值.
\end{enumerate}
\end{problem}

\begin{solution}
\begin{enumerate}
\item 因为点 \(P\) 在椭圆上，所以 \(\frac{1}{5}+\frac{a^2/2}{b^2}=1\)，从而 \(\frac{b^2}{a^2}=\frac58\). 设椭圆焦半距为 \(c\)，则 \(c^2=a^2-b^2=\frac38a^2\)，所以离心率 \(e=\frac{c}{a}=\sqrt{\frac38}=\frac{\sqrt6}{4}\).
\item 由 \(A(-a,0)\)、\(O(0,0)\) 及 \(\lvert AQ\rvert=\lvert AO\rvert=a\)，设 \(Q(x,y)\)，则 \((x+a)^2+y^2=a^2\)，即 \(x^2+y^2+2ax=0\). 点 \(Q\) 不可能为 \(O\)，因为 \(O\) 不在椭圆上；若 \(x=0\)，则由上式还得 \(y=0\)，同样矛盾，所以 \(x\ne0\). 设直线 \(OQ\) 的斜率为 \(k\)，则 \(y=kx\)，代入上式得 \(x=-\frac{2a}{1+k^2}\). 又 \(Q\) 在椭圆上，且 \(b^2=\frac58a^2\)，所以 \(\frac{x^2}{a^2}+\frac{k^2x^2}{b^2}=1\)，即 \(\frac{4}{(1+k^2)^2}\left(1+\frac85k^2\right)=1\). 令 \(u=k^2\)，整理得 \(5u^2-22u-15=0\)，故 \(u=5\) 或 \(u=-\frac35\). 因 \(u\geqslant0\)，所以 \(k=\pm\sqrt5\).
\end{enumerate}
\end{solution}

\begin{problem}
已知函数 \(f(x) = \frac{1}{3} x^3 + \frac{1 - a}{2} x^2 - ax - a\)，\(x \in \R\)，其中 \( a > 0 \).
\begin{enumerate}
\item 求函数 \( f(x) \) 的单调区间；
\item 若函数 \( f(x) \) 在区间 \((-2,0)\) 内恰有两个零点，求 \( a \) 的取值范围；
\item 当 \(a = 1\) 时，设函数 \(f(x)\) 在区间 \([t, t + 3]\) 上的最大值为 \(M(t)\)，最小值为 \(m(t)\)，记 \(g(t) = M(t) - m(t)\)，求函数 \(g(t)\) 在区间 \([-3, -1]\) 上的最小值.
\end{enumerate}
\end{problem}

\begin{solution}
\begin{enumerate}
\item \(f'(x)=x^2+(1-a)x-a=(x+1)(x-a)\). 因为 \(a>0\)，所以 \(-1<a\). 当 \(x<-1\) 或 \(x>a\) 时，\(f'(x)>0\)；当 \(-1<x<a\) 时，\(f'(x)<0\). 因此 \(f(x)\) 的增区间为 \((-\infty,-1)\) 和 \((a,+\infty)\)，减区间为 \((-1,a)\).
\item \(f(-2)=-\frac23-a<0\)，\(f(-1)=\frac{1-3a}{6}\)，\(f(0)=-a<0\). 在 \((-2,-1)\) 上函数递增，在 \((-1,0)\) 上函数递减，因此 \(-1\) 是区间 \((-2,0)\) 内的唯一可能的极大点. 当 \(f(-1)>0\)，即 \(0<a<\frac13\) 时，由连续性和严格单调性，\((-2,-1)\) 与 \((-1,0)\) 内各有且仅有一个零点. 当 \(a=\frac13\) 时，\(f(-1)=0\)，但两侧函数值均小于 \(0\)，所以区间内只有零点 \(x=-1\) 一个；当 \(a>\frac13\) 时，\(f(-1)<0\)，函数在整个 \((-2,0)\) 内均小于 \(0\)，没有零点. 所以 \(a\in\left(0,\frac13\right)\).
\item 当 \(a=1\) 时，\(f(x)=\frac13x^3-x-1\)，且 \(f'(x)=x^2-1\). 在区间 \([t,t+3]\)，其中 \(t\in[-3,-1]\)，函数在 \(x=-1\) 处取得最大值，且 \(f(-1)=-\frac13\)，因此 \(M(t)=-\frac13\).

当 \(-3\leqslant t\leqslant-2\) 时，\(t+3\in[0,1]\)，最小值在两个端点中取得. 又 \(f(t+3)-f(t)=3(t+1)(t+2)\geqslant0\)，所以 \(m(t)=f(t)\)，从而 \(g(t)=-\frac13-f(t)\). 因为 \(f\) 在 \([-3,-2]\) 上递增，故此时 \(g(t)\geqslant g(-2)\).

当 \(-2\leqslant t\leqslant-1\) 时，区间 \([t,t+3]\) 包含 \(x=1\)，且 \(f(1)=-\frac53\). 两端点分别属于 \([-2,-1]\) 和 \([1,2]\)，函数值均不小于 \(-\frac53\)，故 \(m(t)=f(1)=-\frac53\)，从而 \(g(t)=\frac43\). 另外 \(f(-2)=-\frac53\)，所以 \(g(-2)=-\frac13+\frac53=\frac43\). 因此 \(g(t)\) 在 \([-3,-1]\) 上的最小值为 \(\frac43\).
\end{enumerate}
\end{solution}
